\documentclass[11pt,letterpaper]{article}
\usepackage[spanish]{babel}
\usepackage[utf8]{inputenc}
\usepackage[T1]{fontenc}
\usepackage{lmodern}
\usepackage{geometry}
\usepackage{amsmath,amssymb}
\usepackage{enumitem}
\geometry{margin=2.5cm}
\setlength{\parindent}{0pt}
\setlength{\parskip}{7pt}
\begin{document}
\begin{center}
{\Large\bfseries Taller de pruebas de hipótesis}\\[3pt]
{\large Aplicaciones en Administración de Empresas y Mercadeo}
\end{center}


Objetivo: Formular y contrastar hipótesis estadísticas en situaciones de toma de decisiones empresariales. Justifique los supuestos, seleccione la prueba apropiada, calcule el estadístico y el valor p, y redacte una conclusión en términos del problema.

\section*{Instrucciones generales}
Para cada problema: (a) defina los parámetros poblacionales; (b) formule \(H\_0\) y \(H\_1\); (c) establezca el nivel de significancia; (d) verifique los supuestos pertinentes; (e) calcule el estadístico de prueba y el valor p; (f) decida sobre \(H\_0\); (g) interprete el resultado y recomiende una acción empresarial. Cuando compare medias de grupos independientes, evalúe la igualdad de varianzas y elija justificadamente la prueba t apropiada. No confunda ausencia de evidencia para rechazar \(H\_0\) con demostración de que \(H\_0\) es verdadera.

\section*{Problema 1. Tasa De Conversión De Una Campaña Digital}
Una empresa de comercio electrónico considera rentable una campaña si su tasa de conversión supera el 5 \%. Tras rediseñar los anuncios, observa 3 compras entre 40 visitas seleccionadas aleatoriamente. ¿Existe evidencia de que la tasa de conversión es superior al 5 \%? Use alfa = 0,05. Indique si resulta adecuado emplear la aproximación normal o si conviene una prueba binomial exacta.

\section*{Problema 2. Productividad De Dos Equipos Comerciales}
La dirección compara las ventas semanales (en millones de pesos) de dos equipos independientes. Equipo A: n = 36, media = 6,0 y desviación estándar = 4,0. Equipo B: n = 40, media = 8,2 y desviación estándar = 4,3. ¿Hay evidencia de que el equipo B logra mayores ventas promedio? Use alfa = 0,05. Examine la igualdad de varianzas antes de seleccionar el contraste de medias.

\section*{Problema 3. Efectividad De Dos Canales De Publicidad}
Se seleccionan aleatoriamente 12 campañas de Facebook Ads y 12 de TikTok Ads. El retorno sobre inversión publicitaria (ROAS) presenta los siguientes resúmenes: Facebook Ads: media = 3,21, desviación estándar = 5,42; TikTok Ads: media = 4,14, desviación estándar = 3,40. ¿Difieren los ROAS medios de ambos canales? Use alfa = 0,05. Evalúe la homogeneidad de varianzas y explique qué implicaría el resultado para la asignación presupuestal.

\section*{Problema 4. Tiempo De Atención Al Cliente}
Una cadena comercial prueba un nuevo protocolo para reducir el tiempo de atención en caja. En muestras independientes de nueve clientes por procedimiento, se obtienen: protocolo tradicional, media = 35,22 minutos y desviación estándar = 4,94; protocolo nuevo, media = 31,55 minutos y desviación estándar = 4,47. ¿Puede afirmarse que el protocolo tradicional tiene un tiempo promedio mayor que el nuevo? Use alfa = 0,05. Compruebe los supuestos pertinentes y explique si recomendaría implementar el cambio.

\section*{Problema 5. Capacitación Y Desempeño De Vendedores}
Se registran las ventas diarias (en unidades) de 16 vendedores antes y después de un programa de capacitación. Para cada vendedor se calcula la diferencia D = ventas después - ventas antes. Las diferencias presentan media = 2,3125 unidades y desviación estándar = 2,7010 unidades. ¿La capacitación incrementó las ventas promedio? Realice una prueba para datos pareados con alfa = 0,05 y repita la decisión con alfa = 0,01. Explique la importancia de la dependencia entre observaciones.

\section*{Problema 6. Recordación De Marca Entre Dos Empaques}
Una empresa compara la recordación de marca de dos diseños de empaque en grupos independientes de consumidores. Diseño nuevo: n = 61, puntuación media = 40, desviación estándar = 24,9. Diseño anterior: n = 61, puntuación media = 29, desviación estándar = 22,7. ¿Existe una diferencia en la puntuación media de recordación? Use alfa = 0,05. Evalúe la homogeneidad de varianzas y determine si el nuevo diseño ofrece una mejora estadísticamente respaldada.

\section*{Problema 7. Duración De Llamadas De Venta}
Dos equipos de telemercadeo aplican guiones distintos. Se desea saber si el guion del equipo A produce llamadas de menor duración promedio que el del equipo B, sin deteriorar la calidad de atención. Se tienen dos muestras independientes: equipo A, n = 11, media = 7,8 minutos y desviación estándar = 2,1; equipo B, n = 11, media = 10,3 minutos y desviación estándar = 3,1. Use alfa = 0,05. Contraste primero la igualdad de varianzas, después las medias y explique por qué una menor duración no prueba por sí sola una mayor calidad del servicio.

\section*{Problema 8. Proporción De Clientes Que Responden A Una Promoción}
Una empresa prueba dos mensajes promocionales en grupos independientes. Con el mensaje A, 40 de 200 clientes realizan una compra; con el mensaje B, 44 de 200 compran. ¿Existe evidencia de que el mensaje B tiene una tasa de respuesta mayor que el mensaje A? Use alfa = 0,05. Calcule las proporciones muestrales, la proporción combinada bajo \(H\_0\) y el estadístico de contraste. Proponga una decisión sobre la campaña.

\section*{Problema 9. Cambio En El Gasto Promedio Por Cliente}
Una tienda registra el gasto de los mismos diez clientes habituales antes y después de introducir un programa de fidelización. Defina D = gasto después - gasto antes, en miles de pesos. Las diferencias observadas son: 3, 8, 5, 1, 7, 4, 9, 6, 2 y 4. ¿Existe evidencia de que el programa aumentó el gasto promedio por cliente? Use alfa = 0,05. Calcule la media y la desviación estándar de las diferencias y aplique la prueba t pareada. Discuta una limitación de inferir causalidad a partir de un diseño antes-después sin grupo de control.

PREGUNTA DE CIERRE
Compare las decisiones obtenidas en los nueve problemas. Identifique en cuáles se estudian proporciones, medias de grupos independientes y medias pareadas. Explique la diferencia entre significancia estadística e importancia económica o gerencial, y señale al menos dos situaciones en las que se requeriría información adicional antes de adoptar una decisión.

\end{document}